\documentclass[10pt,a4paper]{amsart}
\usepackage[margin=19mm]{geometry}
\usepackage{amsmath,amssymb,mathtools}
\usepackage[scr=rsfs,cal=euler]{mathalfa}
\usepackage{xfrac}
\usepackage[unicode,hidelinks,bookmarks=false]{hyperref}
\newcommand{\ie}{i.e.\ }
\newcommand{\cf}{cf.\ }
\newcommand{\resp}{respectively\ }
\newcommand{\Subsection}{\subsection}
\providecommand{\nobreakdash}{\nobreak\hbox{-}}
\setlength{\emergencystretch}{2em}
\multlinegap0pt
\usepackage{enumerate}
\usepackage{amssymb}
\usepackage{xcolor}
\usepackage{dsfont}
\newtheorem{lemma}{Lemma}
\newtheorem{prop}{Proposition}
\usepackage{cleveref}
\crefname{lemma}{Lemma}{Lemmas}
\crefname{prop}{Proposition}{Propositions}
\newcommand{\EKR}{\mathcal{F}}
\newcommand{\Gauss}[2]{{\begin{bmatrix} #1 \\ #2 \end{bmatrix}_q}}
\newcommand{\OO}{\mathcal{O}}
\DeclareMathOperator{\PG}{PG}
\title{Cocliques in the Kneser graph on $(n-1,n)$-flags of $\PG(2n,q)$}
\author{Philipp Heering}
\date{Source: arXiv:2603.09769v1 (CC BY 4.0)}
\begin{document}
\maketitle

\noindent\emph{Source and licence.} Cocliques in the Kneser graph on $(n-1,n)$-flags of $\PG(2n,q)$. Version arXiv:2603.09769v1 (CC BY 4.0), distributed by arXiv under the Creative Commons Attribution 4.0 International licence, http://creativecommons.org/licenses/by/4.0/, as recorded in the arXiv licence metadata for this exact version and shown on the abstract page https://arxiv.org/abs/2603.09769v1. Full licence notice: \texttt{licenses/arxiv-2603.09769-CC-BY-4.0.txt}.\par\smallskip

\noindent\textit{Editorial scope.} Counted unit: the article's prop prop (source0.tex lines 450--452). The article's complete original proof of it is delivered with it (source0.tex lines 454--466, one \texttt{proof} environment of 13 lines). No mathematics is written, repaired or re-derived: the statement and the proof are the article's own lines, in the article's own notation, and no retained line is substituted or re-flowed.  Retained uncounted as the source-local context the proof needs: lines 338--340; lines 353--355; lines 435--437.  Nothing was omitted from any retained span, so the package carries no omission record.  No retained line contains a citation, so the wrapper carries no bibliography and no bibliography entry.  No retained line contains a figure environment, an image-inclusion command or drawing code, so no image is delivered and the package declares no asset dependency.  Editorial formatting only, in addition: an AMS \texttt{amsart} preamble that restates the article's own declarations and macros used by the retained lines, namely the environments \texttt{lemma}, \texttt{prop} on separate counters, together with the standard packages those lines need.  The retained environments print in this document as Lemma 1, Proposition 1 in order of appearance, whereas the article prints the same items with the numbering of its own full document.  The complete native source of the article as distributed in the arXiv e-print for this exact version is bundled unchanged as \texttt{sources/196086/source0.tex}.
\begin{lemma} \label{L: many yellow pw skew}
     Let $\EKR$ be a maximal coclique of $\Gamma_{2n}$ with $n\geq 3$ and assume that the number of red spaces is $\OO(q^{n^2-2})$. Then $|\EKR|$ is $\OO(q^{n^2+n-2})$, or there are at least $n+1$ yellow $(n-1)$-spaces in $\EKR$ that are pairwise skew, or there are at least $n+1$ yellow $n$-spaces in general position.
\end{lemma}

\begin{lemma} \label{L: few that meet all}
 Let $A_1,\ldots,A_{n+1}$ be pairwise skew $(n-1)$-spaces of $\PG(2n,q)$. The number of $n$-spaces that intersect every $A_i$ for $1\leq i\leq n+1$ is $\OO(q^{n^2-1})$.
\end{lemma}

\begin{lemma} \label{L: yellows meet almost all}
    Let $\EKR$ be a maximal coclique of $\Gamma_{2n}$ with $n\geq 3$ and let $B$ be a yellow $n$-space of $\PG(2n,q)$. The number of flags $(A',B')\in \EKR$ with pairwise distinct $n$-spaces and with $A'\cap B=\emptyset$ is $\OO(q^{n^2-1})$.
\end{lemma}

\begin{prop}
Let $n\geq 3$ and let $\EKR$ be a maximal coclique of $\Gamma_{2n}$ such that the number of red spaces in $\EKR$ is $\OO(q^{n^2-2})$. Then $|\EKR|$ is $\OO(q^{n^2+n-2})$.
\end{prop}

\begin{proof}
 The number of red spaces is $\OO(q^{n^2-2})$ and a red space is contained in $\Gauss{n+1}{1}$ flags of $\EKR$, hence the number of flags in $\EKR$ that contain a red space is $\OO(q^{n^2-2})\cdot \Gauss{n+1}{1}=\OO(q^{n^2+n-2})$. 
   % As we have seen already in the proof of \Cref{L: many yellow pw skew}, the number of flags in $\EKR$ with a red space is $\OO(q^{n^2+n-2})$.
   We show that the number of flags in $\EKR$ with a yellow space is also $\OO(q^{n^2+n-2})$. Assume that this is not the case.

     By \Cref{L: many yellow pw skew}, there are $n+1$ pairwise skew yellow $(n-1)$-spaces, or $n+1$ yellow $n$-spaces in pairwise general position in $\PG(2n,q)$. We assume that $\PG(2n,q)$ contains $n+1$ pairwise skew $(n-1)$-spaces $A_1,\ldots,A_{n+1}$ and show that the number of flags in $\EKR$ is $\OO(q^{n^2+n-2})$, for yellow $n$-spaces dual arguments suffice.
     
     Let $1\leq i\leq n+1$. 
     Any flag $(A,B)\in \EKR$ satisfies $B\cap A_i=\emptyset$ or $B\cap A_i\neq\emptyset$.
     By the dual statement of \Cref{L: yellows meet almost all}, the number of flags $(A,B)\in \EKR$ with distinct $(n-1)$-spaces and with $A_i\cap B=\emptyset$ is at most $\OO(q^{n^2-1})$. As the red spaces are already accounted for, we can safely assume that these $(n-1)$-spaces are all yellow. This yields that the number of flags $(A,B)\in \EKR$ with $A_i\cap B=\emptyset$ is $\OO(q^{n^2-1})\Gauss{n}{1}=\OO(q^{n^2+n-2})$. 
     
     Hence, the number of flags $(A,B)\in \EKR$ with $A_i\cap B=\emptyset$ for some $1\leq i\leq n+1$ is $(n+1)\OO(q^{n^2+n-2})=\OO(q^{n^2+n-2})$. All other flags in $\EKR$ must have an $n$-space that meets $A_1,\ldots,A_{n+1}$. By \Cref{L: few that meet all}, the number of $n$-spaces that meet $A_1,\ldots,A_{n+1}$ is $\OO(q^{n^2-1})$. As the red spaces are already accounted for, we can safely assume that these $n$-spaces are all yellow. This yields that the number of flags $(A,B)\in \EKR$ with $A_i\cap B\neq\emptyset$ for all $1\leq i\leq n+1$  is also $\OO(q^{n^2-1})\Gauss{n}{1}=\OO(q^{n^2+n-2})$. 
\end{proof}

\end{document}
